% concatenated from main.tex
\documentclass{article}
\usepackage{graphicx} % Required for inserting images
\usepackage{amsmath,times,hyperref, multicol}
\usepackage{amsthm}
\usepackage{amssymb}
\usepackage{amsfonts}
\usepackage{latexsym}
\usepackage{graphicx}
\usepackage[usenames,dvipsnames]{color} 
\usepackage{enumitem}
\usepackage{halloweenmath}
\setlist[itemize]{leftmargin=*}
\setlist[enumerate]{leftmargin=*}
\usepackage[font={small}]{caption}
\usepackage{subcaption}
\usepackage{cite}
\usepackage{comment}
\usepackage{mathtools}
\mathtoolsset{showonlyrefs}



\newtheorem{Theorem}{Theorem}[section]
\newtheorem{Lemma}[Theorem]{Lemma}
\newtheorem{Proposition}[Theorem]{Proposition}
\newtheorem{Corollary}[Theorem]{Corollary}
\newtheorem{Definition}[Theorem]{Definition}
\newtheorem{Question}[Theorem]{Question}
\newtheorem{Example}[Theorem]{Example}
\newtheorem{Remark}[Theorem]{Remark}
\newtheorem{Conjecture}[Theorem]{Conjecture}



\title{$q$-Numerical Ranges and Spectral Sets}
\author{
Ryan O'Loughlin\thanks{Department of Mathematics and Statistics, University of Reading,
Reading, England\\ \texttt{r.d.oloughlin@reading.ac.uk}}
\and
Jyoti Rani\thanks{Indian Institute of Science Education and Research (IISER) Mohali, Knowledge City, S.A.S Nagar, Punjab 140306\\ \texttt{jyotirani@iisermohali.ac.in}}
}
\date{}

\begin{document}

\maketitle

\begin{abstract}
We study spectral constants for convex domains $\Omega$ containing the spectrum of an operator. We extend the Crouzeix–Palencia framework by obtaining bounds depending on a parameter $\gamma$ and relating these bounds to geometric properties of $\Omega$ and the numerical range $W(A)$. We generalise the proof that the numerical range is a $1+\sqrt{2}$-spectral set to scaled $q$-numerical ranges. We also propose a generalisation of Crouzeix's Conjecture in the context of $q$-numerical ranges.
\end{abstract}

\medskip

\noindent\textbf{Keywords:} numerical range, spectral set, Crouzeix conjecture, $q$-numerical range, spectral constants

\medskip

\noindent\textbf{MSC (2020):} 47A12, 47A25, 47A60


\section{Introduction}

Let $A$ be a bounded operator on a separable Hilbert space $\mathcal{H}$ and $\Omega \subseteq \mathbb{C}$ be a bounded connected set such that the spectrum is contained in the closure of $ \Omega $. If for all $f \in {\cal A} ( \Omega )$ (holomorphic functions on the interior of $\Omega$ which are continuous on the boundary of $\Omega$) we have
$$
\| f(A) \| \leq c_{\Omega} \sup_{z \in \Omega} |f(z)|
$$
then $\Omega$ is said to be a \textit{$c_{\Omega}$-spectral set}. In this case we call $c_{\Omega}$ the \textit{spectral constant}.

One of the most important results in operator theory is von Neumann's inequality which states that when $\|A\| \leq 1$ the unit disk is a 1-spectral set. Crouzeix's Conjecture claims that the numerical range, $W(A) = \{ \langle Ax,x \rangle :  \|x\| = 1 \}$ is a 2-spectral set. Thus far, the best known spectral constant for the numerical range is $1+ \sqrt{2}$ \cite{1root2, ransford2018remarks}. Crouzeix's Conjecture has gained a lot of traction from the Functional and Complex Analysis communities \cite{o2024crouzeix, CC11.08, CC2matrix, CC3by3, CCgorkin, CCincreasing, CCLi, CCpalenciaextension, CCnilpotent, CCnum, malman2024double, 1root2, CCpalenciaextension, mashreghi2025schwarz}. The main result of this paper provides spectral constant bounds for the scaled
$q$-numerical range $\Omega_q = W_q(A)/q$. These bounds depend on the parameter
$q$ and on a quantity measuring the non-normality of $A$. When $q=1$ the result recovers the
$1+\sqrt{2}$ spectral constant bound for the numerical range obtained in
\cite{1root2}.

In Section \ref{2} we extend previous results in \cite{CCpalenciaextension,crouzeix2019spectral}, and deduce a spectral constant for $\Omega$ which depends on a parameter $\gamma$ coming from a positivity criterion on the double layer potential. In Section \ref{3} we obtain explicit spectral constant bounds for operators $A$ only dependent on the geometry of $\Omega$ and $W(A)$. In Section \ref{4} we specialise this to the case when $\Omega = \Omega_q$ is a scaled $q$-numerical range,
$$
\Omega_q = \frac{W_q(A)}{q}, \quad W_q(A) = \{ \langle Ax,y \rangle : \|x\| = \|y\| = 1, \, \langle x,y \rangle = q \}, 
$$
where $q$ is a complex number such that $ 0< |q| \leq 1$ to deduce spectral constants for $\Omega_q$, which generalises the $1 + \sqrt{2}$ spectral constant result for $ W_1(A)$ in \cite{1root2}. Finally, we conjecture an optimal spectral constant for $\Omega_q$, which generalises the conjectured 2-spectral constant for $W(A)$.















% % One such constant can be derived from the Cauchy integral formula:
% % \begin{equation}
% % \| f(A) \| \leq \frac{1}{2 \pi} \left( \int_{\partial \Omega} \| ( \sigma I - A )^{-1} \|~
% % | d \sigma | \right) \| f \|_{\Omega} , \label{Cauchybound}
% % \end{equation}
% % but this usually is not optimal. 

% The numerical range $W(T)$, numerical radius $\omega(T)$, and operator norm $\|T\|$ of an operator $T \in \mathcal{B(H)}$ are defined as follows: \begin{align*} W(T)&:= \{\langle Tx,x \rangle: x \in \mathbb{S}^1\}, \\
% \omega(T)&:= \sup\{| \langle Tx,x \rangle |: x \in \mathbb{S}^1\} \end{align*} and \begin{align*} \|T\|:= \sup\{| \langle Tx,y \rangle |: x,y \in \mathbb{S}^1\}, \end{align*} where $\mathbb{S}^1$ denotes the unit circle in $\mathcal{H}$, consisting of all vectors $x \in \mathcal{H}$ with $\|x\|=1$.It is generally known that $\omega(\cdot)$ defines a norm on $\mathcal{B(H)}$, which is equivalent to the ordinary operator norm $\|\cdot\|$. In fact, for every $T \in \mathcal{B(H)}$, the following inequalities holds: \begin{equation*}
% \frac{\|T\|}{2} \le w(T) \le \|T\|.
% \end{equation*}

% The classical numerical range has been expanded in numerous directions. Among these extensions, we mention the joint numerical range, the $q$-numerical range, and the $A$-numerical range.
% In the following, consider $q \in \mathbb{C}$ with $|q| \leq 1$ and $\dim \mathcal{H} \geq 2$. Note that the $q$-numerical range $W_q(T)$ and $q$-numerical radius $\omega_q(T)$ of an operator $T \in \mathcal{B(H)}$ are defined as follows:
% \begin{align*}
% W_q(T) &:= \{\langle Tx, y \rangle: (x, y) \in \mathbb{S}_q(\mathcal{H})\},
% \end{align*}
% and
%  \begin{equation*}
% \omega_q(T):= \sup\{|\langle Tx, y \rangle|: (x, y) \in \mathbb{S}_q(\mathcal{H})\}
% \end{equation*}
% % It is well-known that $W_q(T)$ is a non-empty subset of $\mathbb{C}$, see \cite{li1994generalized}. 


% Throughout $A$ will be a bounded operator on a separable Hilbert space, and (unless otherwise stated) $\Omega$ will denote a smoothly bounded convex subset of $\mathbb{C}$ which contains the spectrum of $A$ in its interior. Other domains may be considered, but since $\kappa$ is a spectral constant for $A$ on $X \subseteq \mathbb{C}$ if and only if $\kappa$ is a spectral constant for $A$ on every smoothly bounded neighbourhood $\Omega$ of the spectrum of $A$ for which $\Omega$ contains the closure of $X$ \cite{schwenninger2025double}, we restrict ourself to this setting.
Throughout, $A$ denotes a bounded operator on a separable Hilbert space $\mathcal{H}$. Unless stated otherwise, $\Omega \subseteq \mathbb{C}$ denotes a smoothly bounded convex domain containing the spectrum of $A$ in its interior. 
This restriction entails no loss of generality since $X \subseteq \mathbb{C}$ is a $\kappa$-spectral set if and only if every smoothly bounded neighbourhood $\Omega$ of the spectrum of $A$ containing $\overline{X}$ is also a $\kappa$-spectral set \cite{schwenninger2025double}.

\section{Spectral Constants}\label{2}
Recall ${\cal A} ( \Omega )$ denotes the holomorphic functions on the interior of $\Omega$ which are continuos on the boundary of $\Omega$. For any function $f \in {\cal A} ( \Omega ) $, we define $f(A)$ via the Cauchy integral formula as
\[
f(A) = \frac{1}{2 \pi i} \int_{\partial \Omega} ( \sigma I - A )^{-1} f( \sigma )\,d \sigma .
\]
If $f(z) = \sum_{i=1}^k a_iz^i $ is a polynomial with $a_i \in \mathbb{C}$, then the functional calculus above becomes $f(A) = \sum_{i=1}^k a_iA^i$. If $|\partial \Omega| = L$ and $\sigma:  [0, L] \to \partial \Omega$ is an arc length parametrisation of $\partial \Omega$, then
\[
f(A) = \int_0^L \frac{\sigma' (s)}{2 \pi i} R( \sigma (s) , A ) f( \sigma (s))\,ds ,
\]
where $R( \sigma , A )$ is the resolvent, $( \sigma I - A )^{-1}$.

A key idea in \cite{1root2} was to look at
\begin{equation}\label{cauchytransform}
g(A) := \int_0^L \frac{\sigma' (s)}{2 \pi i} R( \sigma (s) , A ) \overline{f( \sigma (s) )}\,ds,
\end{equation}
which in the scalar case reduces to the Cauchy transform of $\overline{f}$.
Crouzeix and Palencia analysed the operator
\[
 f(A) + g(A )^{*} = \int_0^L \mu ( \sigma (s) , A ) f( \sigma (s) )\,ds , 
\]
where the Hermitian operator $\mu ( \sigma (s) , A )$ is
\begin{equation}
\mu ( \sigma (s) , A ) = \frac{\sigma' (s)}{2 \pi i} R( \sigma (s), A ) + \left[ 
\frac{\sigma' (s)}{2 \pi i} R( \sigma (s), A ) \right]^{*} . \label{mu}
\end{equation}
They showed that $\|f(A) + g(A)^*\| \leq 2$, and then consequently showed that the numerical range was a $1 + \sqrt{2}$-spectral set (see also \cite{ransford2018remarks} for an alternative shorter proof). In \cite{CCpalenciaextension}, the authors studied regions $\Omega$
that do not necessarily contain $W(A)$. They defined 
\begin{equation}
M( \sigma , A ) := \mu ( \sigma , A ) - \lambda_{\min} ( \mu ( \sigma , A ) ) I , \label{M}
\end{equation}
where $\lambda_{\min} ( \mu ( \sigma , A ))$ is the infimum of the spectrum of 
$\mu ( \sigma , A )$ at the point $\sigma$ on $\partial \Omega$.  Thus, by definition, 
$M( \sigma , A )$ is positive semidefinite on $\partial \Omega$.

We introduce the notation 
\begin{equation}\label{gamma}
    \gamma(f) := - \int_0^L \lambda_{\min}( \mu ( \sigma (s), A ) ) 
f( \sigma(s) )\,ds , \quad |\partial \Omega| = L .
\end{equation}

Using the same method of proof as in \cite{1root2}, they proved the following lemma. We provide an alternate proof.

\begin{Lemma}\label{mainlemma}
Let $\Omega$ be a region with smooth boundary containing the spectrum of $A$ in its interior.  
For $f \in {\cal A} ( \Omega )$ with $\| f \|_{\Omega} = 1$, let
\begin{equation}
S := f(A) + g(A )^{*} + \gamma(f) I .
\label{S}
\end{equation}
Then $\| S \| \leq 2 + \gamma(1)$.
\end{Lemma}

\begin{proof}

Let $\Phi(h) = \int_0^L h (s) M(\sigma(s),A) ds$ be a map from continuous function on $\partial \Omega$ to the space of bounded operators on $\mathcal{H}$. Then since $M$ is positive semidefinite, $\Phi$ is a positive map between unital $C^*$ algebras and thus by \cite[Chapter 2]{paulsen2002completely} $\|\Phi\| = \Phi(1) = 2+\gamma(1)$. So for $f \in {\cal A} ( \Omega )$ with $\|f\|_{\Omega} \leq 1$, we have $\|S\| = \|\Phi(f)\| \leq 2 + \gamma(1) $.
\end{proof}
\noindent In \cite{CCpalenciaextension}, the authors then used this lemma to prove the following theorem.

\begin{Theorem}
Let $\Omega$ be bounded, convex domain with smooth boundary containing the spectrum of $A$
in its interior.  Then
$\Omega$ is a $ 1 + \frac{\gamma(1)}{2} + \sqrt{2 + \gamma(1) + \gamma(1)^2 / 4 + m}$-spectral set where 
\begin{equation}
m = \int_0^L | \lambda_{\min} ( \mu ( \sigma (s) , A )) |\,ds .
\label{gamma(1)_hat}
\end{equation}
\label{theorem:1}
\end{Theorem}
\noindent Theorem 2 in \cite{crouzeix2019spectral} sharpens this spectral constant to 
$$
1 + \frac{\gamma(1)}{2} + \sqrt{\left(1+ \frac{\gamma(1)}{2}\right)^2 + c_1 + \hat{\gamma}}, \quad \hat\gamma  :=\max\{|\gamma(f)|\,: |f |\leq 1 \text{ in }\Omega\}
$$
where $c_1 = \sup \{ \|g\|_{\infty, \Omega} : \|f\|_{\infty, \Omega}\leq 1 \}$ and $g$ is the Cauchy transform of $\overline{f}$. We sharpen this result with the following theorem.
\begin{Theorem}\label{main}
    \begin{enumerate}
        \item Let $\Omega$ be bounded, convex domain with smooth boundary containing the spectrum of $A$
in its interior.  Then
$\Omega$ is a \begin{equation}\label{maineq}
    1 + \frac{\gamma(1)}{2} + \sqrt{\left(1+ \frac{\gamma(1)}{2}\right)^2 + a(\Omega)}
\end{equation} 
-spectral set, where $a(\Omega) = \sup_{f \in \mathcal{A}(\Omega)}\{ \inf_{\lambda \in \mathbb{C}}\|g- \lambda \|_{\infty , \Omega} :  \|f\|_{\infty, \Omega} \leq 1  \}$, where $g$ is the Cauchy transform of $\overline{f}$ .
\item There is a bounded, convex $\Omega$ with smooth boundary containing the spectrum of $A$ and a $f \in \mathcal{A}(\Omega)$ such that $\|f(A)\| = \left(1 + \frac{\gamma(1)}{2} + \sqrt{\left(1+ \frac{\gamma(1)}{2}\right)^2 + a(\Omega)}\right)\|f\|_{\Omega}$.
\end{enumerate}
\end{Theorem}
The second statement in the theorem above shows that spectral constant $1 + \frac{\gamma(1)}{2} + \sqrt{\left(1+ \frac{\gamma(1)}{2}\right)^2 + a(\Omega)}$ is in some sense optimal. In \cite{malman2024double} the quantity $a(\Omega)$ was studied to sharpen the $1+ \sqrt{2}$ spectral constant for the numerical range by an indefinite amount. Before proving Theorem \ref{main} we state a result, which will be instrumental in the proof of Theorem \ref{main}.




In the following theorem a function $\hat{f} \in \mathcal{A}(\Omega)$ maximising $\|f(A)\|$ among all holomorphic functions $f$ on $ \Omega$ with $|f| \leq 1$ on $ \Omega$ is called extremal for the pair $(A, \Omega)$. A $x \in \mathcal{H}$ is said to be an associated extremal vector if $\|\hat{f}(A)\|\|x\| = \|\hat{f}(A)x\|$. By construction, if $\|\hat{f}(A)\| \leq c_\Omega$, then $\Omega$ is a $c_{\Omega}$-spectral set.

\begin{Theorem}\label{orthogthm} Let $\hat{f}$ be extremal for $(A,\Omega)$, and let $x$ be an associated extremal unit vector.
Then there exists a unique Borel probability measure $\mu$ on $\partial\Omega$ such that
\begin{equation}\label{1st}
\langle h(A)x,x\rangle=\int_{\partial\Omega} h\,d\mu \quad \text{ for all } h\in \mathcal{A}(\Omega).
\end{equation}
If, further, $\|\hat{f}(A)\| > 1$, then \begin{equation}\label{E:rep}
    \int_{\partial \Omega} \hat{f} d\mu = 0.
\end{equation}
\end{Theorem}


\noindent \textit{Proof of Theorem \ref{main}.}
    \begin{enumerate}
        \item Let $\hat{f}$ be an extremal function with associated extremal unit vector $x$ and let $g$ be the Cauchy transform of $\overline{\hat{f}}$. Then by Lemma \ref{mainlemma} we have 
\begin{align}
    & \|\hat{f}(A)\|^2 + \inf_{\lambda \in \mathbb{C}} \left(\Re \int_{\partial\Omega}\hat{f}(g- \lambda) d \mu \right)  \\
    & = \|\hat{f}(A)\|^2 + \Re \int_{\partial\Omega}\hat{f}g d \mu \quad \text{by \eqref{E:rep}} \\
    &\leq |\|\hat{f}(A)\|^2 + |\langle x, \hat{f}g(A)x \rangle| \quad \text{by \eqref{1st}}\\
    &=|\langle Sx,\hat{f}(A)x \rangle| \quad \text{by Lemma \ref{mainlemma} and \eqref{E:rep}} \\
    &\leq (2 + \gamma(1)) \|\hat{f}(A)\| \quad \text{by Lemma \ref{mainlemma}}.
\end{align}
Labelling $r  = \inf_
{\lambda \in \mathbb{C}} \left(\Re \int_{\partial\Omega}\hat{f}g- \lambda) d \mu \right)$ algebraically manipulating the above yields $$\|\hat{f}(A)\| \leq  1 + \frac{\gamma(1)}{2} + \sqrt{\left(1+ \frac{\gamma(1)}{2}\right)^2 - r}.$$
Since $\mu $ is a probability measure and $\|\hat{f}\|_{\Omega} \leq 1$ it follows that $|r| \leq a(\Omega)$, and thus statement $(i)$ follows. 
\item This statement follows by setting $A = \begin{pmatrix}
    0 & 2 \\
    0 & 0
\end{pmatrix}$, $ \Omega = W(A) = \{z \in \mathbb{C} :  |z| \leq 1 \}$ and $f(z) = z$. In this case, a direct computation yields $\lambda_{\min} ( \mu ( \sigma (s) , A )) = 0$ and $g = \overline{f(0)}$, and thus $a(\Omega) = 0$. 
\end{enumerate}
\qed

\begin{Remark}
Note that when $\Omega = W(A)$ we recover the spectral constant obtained in \cite{malman2024double}.
\end{Remark}


\section{Geometric Bounds for Spectral Constants}\label{3}
From Theorem \ref{main} we see the spectral constant for $\Omega$ depends on the parameter $\gamma(1)$ (where $\gamma$ is defined in \eqref{gamma}). In this section we estimate the parameter $\gamma(1)$ to obtain spectral constants dependent on the geometry of $\Omega$ and $W(A)$. The following two lemmas are needed for the proof of Theorem \ref{maingamma}.

The function,
$$
h_{\Omega}(e^{i \theta}) := \sup_{p \in \Omega} \Re(pe^{-i \theta})
$$
is called the support function of $\Omega$.

\begin{Lemma}\label{per}
$$
\int_0^{2 \pi} h_{\Omega}( e ^{i \theta}) d \theta = |\partial \Omega|.
$$
\end{Lemma}
\begin{proof}
The projection of $\Omega$ on $e^{i \theta} \mathbb{R}$ has length $h_{\Omega}( e ^{i \theta})  + h_{\Omega}( -e ^{i \theta})$, so Cauchy's surface area formula (see, for example Theorem 5.2.1 in \cite{Schneider2014}) yields \begin{equation}\label{perimetersupport}
        \int_0^{2 \pi}  h_{\Omega}(e^{ i \theta})  d \theta = |\partial \Omega|.
    \end{equation}
\end{proof}

\begin{Lemma}\label{3.1}
\begin{enumerate}
    \item Let $ W(A) \subseteq \Omega$ and suppose $W(A)$ has non-empty interior and $v$ be the unit normal at $\sigma \in \partial \Omega$. Then 
\begin{equation}\label{lambdaminforqnumrange}
      \frac{1}{\pi} r(\sigma) \left( h_{\Omega}(v) - h_{W(A)}(v) \right) \leq \lambda_{\min} ( \mu ( \sigma  , A )),
\end{equation}
where $r(\sigma) = \inf_{\|y\|=1}\|(\sigma - A)^{-1}y\|^2$ is the square of the minimum modulus of the resolvent.
\item     Let $w_{\Omega}(A)= \sup_{z \in \Omega}(|z|)$ and $w_{W(A)} = \sup_{z \in W(A)}(|z|)$. Then 
    $$
    \left(\frac{1}{w_{\Omega}(A) + 2w_{W(A)}}\right)^2 \leq r(\sigma).
    $$
\end{enumerate}    
\end{Lemma}
\begin{proof}
\begin{enumerate}
    \item Let $\|y\| =1 $ and $y = (\sigma - A)x$. As $ \mu ( \sigma (s) , A )$ is a selfadjoint matrix, we have
    \begin{align}
        &\lambda_{\min} ( \mu ( \sigma (s) , A )) \\
        &= \inf_{\|y\|=1}\langle\mu ( \sigma (s) , A )y,y \rangle \\
        & =\inf_{\|y\| = 1} \left( \frac{1}{\pi} \|x\|^2\Re\left(\left(\sigma - \left\langle A \frac{x}{\|x\|}, \frac{x}{\|x\|} \right)\right\rangle\right)v^*\right) \\
        & \geq \frac{1}{\pi} r(\sigma) \inf_{p \in W(A)} \Re ((\sigma - p)v^*) \\
        &= \frac{1}{\pi} r(\sigma) \left( \Re(\sigma v^*) - h_{W(A)}(v) \right) \\
        &=\frac{1}{\pi} r(\sigma) \left( h_{\Omega}(v) - h_{W(A)}(v) \right),
    \end{align}
    where $\Re(\sigma v^*) =  h_{\Omega}(v) $ because convexity of $\Omega$ and the supporting half plane condition yields $\Re (zv^*) \leq \Re(\sigma v^*)$ for all $z \in \Omega$.
    \item Observe for each $y$ of unit norm, we have $y = (\sigma I - A)(\sigma I - A)^{-1}y$, so the sub-multiplicative property of norms gives
        $$
        \frac{1}{|\sigma| + \|A\|}\leq \frac{1}{\|(\sigma I - A)\|} \leq \| (\sigma I - A)^{-1}y\|.
        $$
        The result follows since $\|A\| \leq 2w_{W(A)} $.
\end{enumerate}    
\end{proof}






We now present the main result of this section which provides a bound for $\gamma(1)$ only in terms of the geometry of $\Omega$ and $W(A)$. In the following theorem $w_{\Omega}(A)= \sup_{z \in \Omega}(|z|)$ and $w_{W(A)} = \sup_{z \in W(A)}(|z|)$.

\begin{Theorem}\label{maingamma}
Let $W(A)$ have non-empty interior. Then for $\gamma$ as given in \eqref{gamma},
    $$
    \gamma(1) \leq - \frac{1}{\pi }\left(\frac{1}{w_{\Omega}(A) + 2w_{W(A)}}\right)^2 \min_{\sigma \in \partial \Omega}\rho (\sigma)   \left( |\partial \Omega| -  | \partial W(A)| \right),
    $$
    where $\rho (\sigma) $ is the radius of curvature at a point $ \sigma \in \partial \Omega$.
\end{Theorem}

\begin{proof}
    The previous lemma gives
    $$
    \gamma(1) \leq  - \frac{1}{\pi }\left(\frac{1}{w_{\Omega}(A) + 2w_{W(A)}}\right)^2  \int_0^L \left( h_{\Omega}(v) - h_{W(A)}(v) \right) ds .
    $$
    Write $v(s) = e^{i \theta(s)}$, then Frenet formulas imply $\kappa(s) = |v'(s)| = \theta'(s)$ where $\kappa$ is the (unsigned) curvature, so $ds = \rho (\theta) d \theta$ where $\rho(\theta) = \frac{1}{k(\theta)}$ is the radius of curvature. Thus 
    $$
    \gamma(1) \leq  - \frac{1}{\pi }\left(\frac{1}{w_{\Omega}(A) + 2w_{W(A)}}\right)^2   \min_{\sigma \in \partial \Omega}\rho (\sigma) \int_0^{2 \pi} \left( h_{\Omega}(e^{ i \theta}) - h_{W(A)}(e^{ i \theta}) \right) d \theta .
    $$
    Now the result follows from Lemma \ref{per}.
\end{proof}

There are several proofs in the literature (see for example \cite{1root2,okubo1975constants}) showing that if $\Omega$ is a disc with $W(A)\subseteq \Omega$, then $\Omega$ is a $2$-spectral set for $A$. The result above, combined with Theorem~\ref{main}, yields an explicit improvement in the case when $W(A)\subsetneq \Omega$. It provides a computable constant $K<2$, depending only on the geometry of $W(A)$ and $\Omega$, such that $\Omega$ is a $K$-spectral set for $A$.







\section{Scaled q-Numerical Range as a Spectral Set}\label{4}




In this section for $A \in \mathcal{L}(\mathcal{H})$ such that $A \neq \lambda I$ for any $\lambda \in \mathbb{C}$ we consider the scaled $q$-numerical range $\Omega_q:=\frac{W_q(A)}{q}$ for $0<|q| < 1$ as a spectral set. In general $W_q(A)$ cannot be a spectral set since there are matrices with eigenvalues lying outside $W_q(A)$. By contrast, $W(A)\subseteq \Omega_q$ \cite[p. 381]{gau2021numerical}, and since $W(A)$ is a spectral set $\Omega_q$ will also be a spectral set for $A$. Note that $W_q(A)$ (and hence $\Omega_q$) is convex \cite[p. 381]{gau2021numerical} and the spectrum of $A$ lies in the interior of $\Omega_q$ \cite[Theorem 2.10]{li1994generalized}. Throughout this section, in order to apply the results of previous sections, (where it was assumed $\Omega$ had smooth boundary) we further assume $A$ is such that $\Omega_q$ has smooth boundary, this holds for example when $\mathcal{H}$ is finite dimensional \cite{rajic2005generalized}.




By applying similar arguments to those used in 
\cite[Theorem 2.5]{li1994generalized} 
in the operator setting, it follows that 
$\Omega_{q_1} \subseteq \Omega_{q_2}$ 
whenever $0 < |q_2| \le |q_1| < 1$, it is natural to expect the corresponding spectral constant for $\Omega_q$ to decrease as $|q|$ decreases. 
% 
Furthermore any nontrivial bound must depend not only on $q$ but also on a quantitative measure of the non-normality of $A$. When $\mathcal{H}$ is finite dimensional, the term
\[
\int_0^{2\pi}\sup_{x\in M_\theta}\sqrt{\|Ax\|^2-|\langle Ax,x\rangle|^2}\,d\theta
\]
 where $M_\theta=\arg\max_{\|x\|=1}\Re\langle e^{-i \theta}A x,x\rangle$ plays this role. However $M_\theta$ may be empty when $\mathcal{H}$ is infinite dimensional, so we consider the set of approximate maximisers $M_{\theta, \varepsilon}$ in the theorem below. Moreover, the furthermore part identifies this integral as a variational quantity, describing the initial rate of change of $|\partial \Omega_q|$ as $|q|$ varies. The following theorem generalises the result that the numerical range is a
$1+\sqrt{2}$-spectral set \cite{1root2} by providing a bound for the spectral constant of
the scaled $q$-numerical range $\Omega_q$. The bound depends on $q$ and on a
quantity measuring the non-normality of $A$.






\begin{Theorem}
For $0 < |q| < 1$, let $\Omega_q = \frac{W_q(A)}{q}$ have smooth boundary and let $W(A)$ have non-empty interior. For $\theta\in[0,2\pi]$ and $\varepsilon>0$, define
\[
M_{\theta,\varepsilon}
:=
\left\{
x\in H:\|x\|=1,\ 
\Re\langle e^{-i\theta}Ax,x\rangle
\ge h_{W(A)}(e^{i\theta})-\varepsilon
\right\},
\]
and set
\[
m(\theta)
:=
\lim_{\varepsilon\downarrow 0}
\sup_{x\in M_{\theta,\varepsilon}}
\sqrt{\|Ax\|^2-|\langle Ax,x\rangle|^2}.
\]
Then
\[
\gamma(1)\le
-C\min_{\sigma\in\partial\Omega_q}\rho(\sigma)
\frac{\sqrt{1-|q|^2}}{|q|}
\int_0^{2\pi} m(\theta)\,d\theta,
\]
where
\[
C=\frac1\pi\left(\frac{1}{w_{\Omega_q}(A)+2w_{W(A)}}\right)^2.
\]

















Furthermore setting $t \in (0, \infty)$ such that $|q| = \frac{1}{\sqrt{1+t^2}}$ and writing $ \Omega(t) =  \Omega_{q(t)}$ we have 
$$
\int_0^{2 \pi} m(\theta) = \frac{d}{dt}_{|_{t=0+}}\left|\partial\Omega(t) \right|.
$$
\end{Theorem}

\begin{proof}
First observe that $\Omega_q = \Omega_{|q|}$ so it suffices to prove the theorem for $\Omega_{|q|}$. Furthermore, in light of the previous theorem and Lemma \ref{per} it suffices to prove that 
$$
\int_0^{2 \pi} \left( h_{\Omega_{|q|}}(e^{ i \theta}) - h_{W(A)}(e^{ i \theta}) \right) ds   \geq \frac{\sqrt{1-|q|^2}}{|q|} \int_0^{2 \pi} m(\theta).
$$
Fix $\theta\in[0,2\pi]$ and $x \in \mathcal{H}$ with $\|x\|=1$. Then any $ y \in \mathcal{H}$ with $\|y\|=1$ and $\langle x,y\rangle=|q|$ can be written as
$y=|q|x+\sqrt{1-|q|^2}\,z$ with $z\perp x$ and $\|z\|=1$. Hence
\begin{align}
    &\Re\big(e^{-i\theta}\langle Ax,y\rangle\big)
= |q|\,\Re\langle e^{-i \theta}A x,x\rangle
+\sqrt{1-|q|^2}\,\Re\langle e^{-i \theta}A x,z\rangle \\
& \le |q|\,\Re\langle e^{-i \theta}A x,x\rangle + \sqrt{1-|q|^2}\,\|P_{x^\perp}(e^{-i \theta}A x)\|,
\end{align}
where $P_{x^\perp}$ is the projection onto the orthogonal complement of the span of $x$. Clearly equality is attained by choosing $z$ in the direction of $P_{x^\perp}(e^{-i \theta}A x)$ and so
\[
h_{W_{|q|}(A)}(e^{i\theta})
=\sup_{\|x\|=1}\Big(|q|\,\Re\langle e^{-i \theta}A x,x\rangle+\sqrt{1-|q|^2}\,\|P_{x^\perp}(e^{-i \theta}A x)\|\Big),
\]
which means
\begin{equation}\label{star1}
    h_{\Omega_{|q|}}(e^{i\theta})
=\sup_{\|x\|=1}\Big(\Re\langle e^{-i \theta}A x,x\rangle+\frac{\sqrt{1-|q|^2}}{|q|}\,\|P_{x^\perp}(e^{-i \theta}A x)\|\Big).
\end{equation}
For every $\varepsilon>0$, restricting the supremum to $x\in M_{\theta,\varepsilon}$ gives
\[
h_{\Omega_{|q|}}(e^{i\theta})
\ge
\sup_{x\in M_{\theta,\varepsilon}}
\left(
\Re\langle e^{-i\theta}Ax,x\rangle
+
\frac{\sqrt{1-|q|^2}}{|q|}
\|P_{x^\perp}(e^{-i\theta}Ax)\|
\right).
\]
Hence
\begin{align}
    &h_{\Omega_{|q|}}(e^{i\theta})-h_{W(A)}(e^{i\theta}) \\
&\ge
\sup_{x\in M_{\theta,\varepsilon}}
\left(
\Re\langle e^{-i\theta}Ax,x\rangle-h_{W(A)}(e^{i\theta})
+
\frac{\sqrt{1-|q|^2}}{|q|}
\|P_{x^\perp}(e^{-i\theta}Ax)\|
\right).
\end{align}
Since
\[
\Re\langle e^{-i\theta}Ax,x\rangle \ge h_{W(A)}(e^{i\theta})-\varepsilon
\qquad \text{for } x\in M_{\theta,\varepsilon},
\]
it follows that
\[
h_{\Omega_{|q|}}(e^{i\theta})-h_{W(A)}(e^{i\theta})
\ge
-\varepsilon+
\frac{\sqrt{1-|q|^2}}{|q|}
\sup_{x\in M_{\theta,\varepsilon}}
\|P_{x^\perp}(e^{-i\theta}Ax)\|.
\]
Using
\[
\|P_{x^\perp}(e^{-i\theta}Ax)\|^2
=
\|e^{-i\theta}Ax\|^2-|\langle e^{-i\theta}Ax,x\rangle|^2
=
\|Ax\|^2-|\langle Ax,x\rangle|^2,
\]
we obtain
\[
h_{\Omega_{|q|}}(e^{i\theta})-h_{W(A)}(e^{i\theta})
\ge
-\varepsilon+
\frac{\sqrt{1-|q|^2}}{|q|}
\sup_{x\in M_{\theta,\varepsilon}}
\sqrt{\|Ax\|^2-|\langle Ax,x\rangle|^2}.
\]
Letting $\varepsilon\downarrow 0$ yields
\[
h_{\Omega_{|q|}}(e^{i\theta})-h_{W(A)}(e^{i\theta})
\ge
\frac{\sqrt{1-|q|^2}}{|q|}\,m(\theta).
\]
Integrating over $[0,2\pi]$, it follows that
\[
\int_0^{2\pi}
\bigl(h_{\Omega_{|q|}}(e^{i\theta})-h_{W(A)}(e^{i\theta})\bigr)\,d\theta
\ge
\frac{\sqrt{1-|q|^2}}{|q|}
\int_0^{2\pi} m(\theta)\,d\theta.
\]
The desired bound for $\gamma(1)$ now follows from Theorem \ref{maingamma} and Lemma \ref{per}.











We now prove the furthermore statement. From \eqref{star1}, noting that $
t=\frac{\sqrt{1-|q|^2}}{|q|}$, 
we deduce
\begin{align}
&\frac{h_{\Omega(t)}(e^{i\theta})-h_{W(A)}(e^{i\theta})}{t} \\
&=
\sup_{\|x\|=1}
\left(
\frac{\Re\langle e^{-i\theta}Ax,x\rangle-h_{W(A)}(e^{i\theta})}{t}
+
\|P_{x^\perp}(e^{-i\theta}Ax)\|
\right).
\end{align}
Fix $\varepsilon>0$. If $x\notin M_{\theta,\varepsilon}$, then $
\Re\langle e^{-i\theta}Ax,x\rangle-h_{W(A)}(e^{i\theta})<-\varepsilon$, 
and hence
\[
\frac{\Re\langle e^{-i\theta}Ax,x\rangle-h_{W(A)}(e^{i\theta})}{t}
+
\|P_{x^\perp}(e^{-i\theta}Ax)\|
<
-\frac{\varepsilon}{t}+\|A\|.
\]
Thus, for sufficiently small $t$, the supremum may be restricted to vectors
$x\in M_{\theta,\varepsilon}$, and so
\[
\frac{h_{\Omega(t)}(e^{i\theta})-h_{W(A)}(e^{i\theta})}{t}
\le
\sup_{x\in M_{\theta,\varepsilon}}
\|P_{x^\perp}(e^{-i\theta}Ax)\|.
\]
Taking $\limsup$ as $t\downarrow0$ gives
\[
\limsup_{t\downarrow0}
\frac{h_{\Omega(t)}(e^{i\theta})-h_{W(A)}(e^{i\theta})}{t}
\le
\sup_{x\in M_{\theta,\varepsilon}}
\|P_{x^\perp}(e^{-i\theta}Ax)\|.
\]
Since this holds for every $\varepsilon>0$, it follows that
\begin{equation}\label{eq1}
    \limsup_{t\downarrow0}
\frac{h_{\Omega(t)}(e^{i\theta})-h_{W(A)}(e^{i\theta})}{t}
\le
m(\theta).
\end{equation}

Let $\delta>0$. If $x\in M_{\theta,\delta t}$, then $
\Re\langle e^{-i\theta}Ax,x\rangle-h_{W(A)}(e^{i\theta})\ge -\delta t$ and so
\[
\frac{h_{\Omega(t)}(e^{i\theta})-h_{W(A)}(e^{i\theta})}{t}
\ge
-\delta+\|P_{x^\perp}(e^{-i\theta}Ax)\|.
\]
Taking the supremum over $x\in M_{\theta,\delta t}$ yields
\[
\frac{h_{\Omega(t)}(e^{i\theta})-h_{W(A)}(e^{i\theta})}{t}
\ge
-\delta+
\sup_{x\in M_{\theta,\delta t}}
\|P_{x^\perp}(e^{-i\theta}Ax)\|.
\]
Taking $\liminf$ as $t\downarrow0$, we obtain
\[
\liminf_{t\downarrow0}
\frac{h_{\Omega(t)}(e^{i\theta})-h_{W(A)}(e^{i\theta})}{t}
\ge
-\delta+m(\theta).
\]
Since $\delta>0$ is arbitrary, it follows that
\begin{equation}\label{eq2}
    \liminf_{t\downarrow0}
\frac{h_{\Omega(t)}(e^{i\theta})-h_{W(A)}(e^{i\theta})}{t}
\ge
m(\theta).
\end{equation}
Combining \eqref{eq1} and \eqref{eq2} gives
\[
\frac{d}{dt}\bigg|_{t=0+} h_{\Omega(t)}(e^{i\theta})=m(\theta).
\]

Observe that
\[
\|P_{x^\perp}(e^{-i\theta}Ax)\|^2
=
\|e^{-i\theta}Ax\|^2-|\langle e^{-i\theta}Ax,x\rangle|^2
=
\|Ax\|^2-|\langle Ax,x\rangle|^2,
\]
so
\[
0\le
\frac{h_{\Omega(t)}(e^{i\theta})-h_{W(A)}(e^{i\theta})}{t}
\le
\sup_{\|x\|=1}\|P_{x^\perp}(e^{-i\theta}Ax)\|
\le \|A\|.
\]
Thus the dominated convergence theorem allows us to integrate the derivative identity,
and by Lemma \ref{per} we obtain
\[
\int_0^{2\pi} m(\theta)\,d\theta
=
\frac{d}{dt}\bigg|_{t=0+} |\partial\Omega(t)|.
\]
\end{proof}
We remark that when $\mathcal{H}$ is finite dimensional, the $m(\theta)$ in the theorem above can be replaced with $\sup_{x\in M_\theta}\sqrt{\|Ax\|^2-|\langle Ax,x\rangle|^2}$. Furthermore when $q=1$ the set $\Omega_q$ reduces to the numerical range $W(A)$ and the
result recovers the $1+\sqrt{2}$ spectral constant bound.


The following conjecture generalises Crouzeix's Conjecture from the classical numerical range to (scaled) $q$-numerical ranges.

\begin{Conjecture}
    Let $A$ be a $n \times n$ matrix, $0 < |q| \leq 1$ and $p$ a polynomial with complex coefficients. Then 
    $$
    \|p(A)\| \leq \max\!\left( 1,\,
\frac{2|q|}{\,1+\sqrt{1-|q|^{2}}\,} \right) \sup_{z \in \Omega_q}|p(z)|.
    $$
\end{Conjecture}
The term $1$ in the maximum arises because a spectral constant is always at least 1, which follows by testing with the constant polynomial $p(z)=1$.





























\newpage

\bibliographystyle{plain}
\bibliography{bibliography.bib}
\end{document}
